۸۳  
۸۲  
ساشا بدترین عملکرد را داشت. او در استراتژی‌های اولیش ناموفق بود، به تعمیراتش زل زد. فقط اجتماعی شدن برایش چالشی بود: او مانند یک جوجه غاز بزرگ با جوش صورت به نظر می‌رسید و حتی کمترین میزان اعتماد به نفسی نداشت. دختران فقط از روی ادب به حرف‌هایش گوش می‌دادند.  

این‌بار نوبت من بود که دور میز بچرخم و بپرسم: «نمره‌ات چنده؟» و «نقاط قوتت چیه؟» و «با چند تا دختر دوست داری بخوابی؟»

صاحب‌نظر عجیب، که بیست سال داشت، با دو زن بوده است. «من جرات نزدیک شدن را دارم و بعضی از زنان را در گذشته جذب کرده‌ام» او گفت و چپ دستش را با آرامش روی صندلی کناری گذاشت. «اما نقطه ضعفم مرحله جذب است. حتی وقتی احساس می‌کنم که آنها به من جذب می‌شوند، هنوز به نتیجه نمی‌رسم.»

جری، که سی و سه سال داشت، با سه زن بوده است. «می‌توانم در کافی‌شاپ‌ها و اکثر محیط‌های کم‌صدا کار کنم، اما در کلوب‌ها راحت نیستم.»

و ساشا، که بیست و دو سال داشت، گفت که با یک زن بوده، هرچند ما مشکوک بودیم که یکی را اغراق کرده است. «من به این بازی مشغولم چون شبیه به بازی Dungeons and Dragons است. وقتی یک نقشه معکوس یا یک روتین یاد می‌گیرم، مثل این است که یک طلسم جدید یا عصا دارم که نمی‌توانم صبر کنم تا از آن استفاده کنم.»

یکی پس از دیگری، ترس‌هایشان و ضبط‌کننده‌های صدایشان را روی میز گذاشتند. کار من این بود که آن‌ها را وارد بازی کنم. باید آنچه را که در سرم بود، به ذهن آن‌ها منتقل می‌کردم.

بخش تدریس کارگاه آسان بود. تمام کاری که باید انجام می‌دادم این بود که مرموز را در مسیر نگه دارم—او صدای خودش را دوست داشت—و به آن‌ها مطالب بدهم. چالش، بخش نمایش خواهی بود.

هنگام صحبت، پسران را به مأموریت‌هایی به سمت میزهای مختلف فرستادیم. آن‌ها را ترغیب کردیم دست به کار شوند، به زبان بدن آن‌ها و واکنش‌های زن‌ها نگاه کردیم، سپس به آن‌ها بازخورد دادیم:

«شما به نمایش تکیه کرده بودید، که نیازمندی نشان می‌داد. صاف بایستید و روی پای عقب‌تان سنگینی کنید، گویی ممکن است هر لحظه دور شوید.»

«شما آن‌ها را با ایستادن بیش از حد طولانی بالای سرشان ناراحت می‌کردید. باید می‌نشستید و محدودیت زمانی برای خود تعیین می‌کردید. مثلاً بگویید: «فقط چند دقیقه می‌توانم بمانم چون باید به دوستانم بپیوندم». به این ترتیب آن‌ها نگران نخواهند بود که تمام شب آنجا خواهید ماند.»

مرموز آن را پایداری اجتماعی دینامیک می‌نامد. ما دائماً توسط، از یک سو، تمایل شدیدمان به داشتن رابطه جنسی با یک دختر و، از سوی دیگر، نیاز به محافظت از خود در هنگام نزدیک شدن، به این سو و آن سو می‌زنیم. او می‌گوید دلیل وجود این ترس، آن است که ما به‌طور تکاملی برای یک زندگی قبیله‌ای سازمان یافته‌ایم، جایی که همه افراد جامعه می‌دانند وقتی مردی رد می‌شود.